\vspace*{\fill} % Vertically center the page content
\section*{Περίληψη}
\label{AbstractGreek}

Η συγκεκριμένη εργασία ασχολείται με το πρόβλημα της αναγνώρισης της ανθρώπινης ομιλίας αξιοποιώντας οπτικά μέσα, δηλαδή βίντεο της ομιλίας. Είναι πρόβλημα συγγενικό με το κλασικό πρόβλημα της αυτόματης αναγνώρισης ομιλίας (Automatic Speech Recognition - ASR) και αποκαλείται συχνά Οπτική Αναγνώριση Ομιλίας (Visual Speech Recognition - VSR) ή Οπτικοακουστική Αναγνώριση Ομιλίας (Audio-Visual Speech Recognition - AVSR), σε περίπτωση που χρησιμοποιούνται και ηχητικά στοιχεία, πέραν των οπτικών. Ο τομέας της ``αναγνώρισης χειλιών'' (lip reading) έχει γνωρίσει τον τελευταίο καιρό σημαντική ανάπτυξη, ιδιαίτερα με την έλευση των Transformers και τη δυνατότητα εκπαίδευσης μοντέλων αξιοποιώντας τη βαθιά μάθηση. Πρόσφατες μελέτες τείνουν να διαχωρίσουν το πρόβλημα της αναγνώρισης της ομιλίας σε 2 υπο-προβλήματα: αφενός την αναγνώριση των φωνημάτων (των ελάχιστων φωνητικών τμημάτων της ανθρώπινης ομιλίας) και αφετέρου τον συνδυασμό των προβλεπόμενων φωνημάτων για την πρόβλεψη λέξεων και, τελικά, φράσεων.

Στα πλαίσια της εργασίας μελετάται η επίλυση του πρώτου υπο-προβλήματος, δηλαδή της πρόβλεψης φωνημάτων από βίντεο ενός ομιλητή. Αναπτύσσεται μοντέλο τεχνητής νοημοσύνης το οποίο δέχεται ως είσοδο βουβό βίντεο της περιοχής των χειλιών και επιστρέφει ως έξοδο τις προβλέψεις του για το ποιο φώνημα εκφέρεται σε κάθε frame του βίντεο. Το μοντέλο αυτό αξιοποιεί ένα 3D συνελικτικό νευρωνικό δίκτυο (CNN) για την εξαγωγή features από το βίντεο και BiLSTM ως Sequence Encoder. Παράλληλα, σχεδιάζεται μοντέλο που δέχεται ως είσοδο τόσο το βίντεο της περιοχής των χειλιών όσο και ηχητικά στοιχεία. Το οπτικο-ακουστικό μοντέλο ακολουθεί λογική παρόμοια με αυτή του οπτικού μοντέλου, αξιοποιώντας CNNs για την εξαγωγή των οπτικών και ακουστικών χαρακτηριστικών. Στα χαρακτηριστικά εξαγόμενα από κάθε τροπικότητα (modality) εφαρμόζεται μηχανισμός cross-attention, ενώ τα αποτελέσματα αυτού συνενώνονται με βάση την αξιοπιστία της κάθε τροπικότητας για κάθε χρονική στιγμή. Τέλος, εφαρμόζεται και πάλι BiLSTM ως Sequence Encoder.

Η πειραματική διαδικασία της εργασίας αυτής αποτελείται από 3 πειράματα. Κατά το πρώτο πείραμα, δοκιμάζεται η επίδοση του οπτικού μοντέλου σε περιβάλλον περιορισμένο τόσο ως προς το λεξιλόγιο όσο και ως προς το συντακτικό του. Στο πρώτο υπο-πείραμα, πραγματοποιείται μία μελέτη αφαίρεσης, από την οποία προκύπτει η προτεινόμενη μεθοδολογία ως αυτή με τις βέλτιστες επιδόσεις. Στο δεύτερο υπο-πείραμα, μελετάται σε μεγαλύτερο βάθος η επίδοση του προτεινόμενου μοντέλου, αξιοποιώντας και αλγοριθμικό μηχανισμό πρόβλεψης λέξεων από τα φωνήματα. Στο δεύτερο πείραμα μελετήθηκε η επίδοση της προτεινόμενης οπτικο-ακουστικής μεθοδολογίας σε περιβάλλον περιορισμένο αλλά επηρεασμένο από θόρυβο, τόσο ηχητικό όσο και οπτικό. Στο πρώτο υπο-πείραμα μελετάται η επίδραση της αύξησης του SNR του ηχητικού σήματος στην επίδοση του μοντέλου, ενώ στο δεύτερο υπο-πείραμα παρουσιάζονται τα αποτελέσματα της αύξησης του οπτικού θορύβου στις προβλέψεις του μοντέλου. Τέλος, το τρίτο πείραμα συγκρίνει το προτεινόμενο οπτικοακουστικό μοντέλο με 4 state of the art μοντέλα που αξιοποιούν μία τροπικότητα το κάθε ένα, σε περιβάλλον με θόρυβο στον ήχο και παραμόρφωση στο βίντεο.

\vspace*{\fill}
